\chapter{原子的基本状况}
\problem{1.1}{若卢瑟福散射用的 $\alpha$ 粒子是放射性物质镭 $\mathrm C'$ 放射的，其动能为 $7.68\times10^6$ 电子伏特。散射物质是原子序数 $Z=79$ 的金箔。试问散射角 $\theta=150^\circ$ 所对应的瞄准距离 $b$ 多大？}
\answerarea
\problem{1.2}{已知散射角为 $\theta$ 的 $\alpha$ 粒子与散射核的最近距离为
\[
r_m=\left(\frac{1}{4\pi\varepsilon_0}\right)\frac{2ZZ'e^2}{Mv^2}
\left(1+\frac{1}{\sin(\theta/2)}\right).
\]
试问上题 $\alpha$ 粒子与散射的金原子核之间的最近距离 $r_m$ 多大？}
\answerarea\clearpage
\problem{1.3}{若用动能为 1 兆电子伏特的质子射向金箔，问质子与金箔原子核可能达到的最小距离多大？又问如用同样能量的氘核（氘核带一个 $+e$ 电荷而质量是质子的两倍，是氢的一种同位素的原子核）代替质子，其与金箔原子核的最小距离多大？}
\answerarea
\problem{1.4}{钍放射的一种 $\alpha$ 粒子的速度为 $1.597\times10^7\,\mathrm{m/s}$，正面垂直入射于厚度为 $10^{-7}\,\mathrm m$、密度为 $1.932\times10^4\,\mathrm{kg/m^3}$ 的金箔。试求所有散射在 $\theta>90^\circ$ 的 $\alpha$ 粒子占全部入射粒子的百分比。已知金的原子量为 197。}
\answerarea\clearpage
\problem{1.5}{$\alpha$ 粒子散射实验的数据在散射角很小（$\theta\leq15^\circ$）时与理论值差得较远，是什么原因？}
\answerarea
\problem{1.6}{已知 $\alpha$ 粒子质量比电子质量大 7300 倍。试利用中性粒子弹性碰撞证明：$\alpha$ 粒子散射“受电子的影响是微不足道的”。}
\answerarea\clearpage
\problem{1.7}{能量为 3.5 兆电子伏特的细 $\alpha$ 粒子束射到单位面积上质量为 $1.05\times10^{-2}\,\mathrm{kg/m^2}$ 的银箔上（习题图 1.1），$\alpha$ 粒子与银箔表面成 $60^\circ$ 角。在离 $\alpha$ 入射线成 $\theta=20^\circ$ 的方向上，离银箔散射区距离 $L=0.12\,\mathrm m$ 处放一窗口面积为 $6.0\times10^{-5}\,\mathrm{m^2}$ 的计数器。测得散射进此窗口的 $\alpha$ 粒子是全部入射 $\alpha$ 粒子的百万分之 29。若已知银的原子量为 107.9，试求银的核电荷数 $Z$。
\fig{1_1}{1.1}}
\answerarea
\problem{1.8}{设想铅（$Z=82$）原子的正电荷不是集中在很小的核上，而是均匀分布在半径约为 $10^{-10}\,\mathrm m$ 的球形原子内。如果有能量为 $10^6$ 电子伏特的 $\alpha$ 粒子射向这样一个“原子”，试通过计算论证这样的 $\alpha$ 粒子不可能被具有上述设想结构的原子产生散射角大于 $90^\circ$ 的散射。这个结论与卢瑟福实验结果差得很远，这就说明原子的汤姆逊模型是不能成立的（原子中电子的影响可以忽略）。}
\answerarea\clearpage
